% Section 2.2, from lectures/L02/appendix/b_classes2.md.

\section{Advanced Discussion on Classes}
\label{c2:sec:classes2}\label{c2:app:b}

\subsection{The keyword \code{explicit}}
\label{c2:sec:explicit}
Consider the simple class \code{Led} below:

\begin{cppcode}
#include <cstdint>

class Led final
{
public:
    Led(const std::uint8_t pin) noexcept
        : myPin{pin}
        , myState{false}
    {}

    void setEnabled(const bool state) noexcept { myState = state; }
    [[nodiscard]] bool isEnabled() const noexcept { return myState; }

private:
    const std::uint8_t myPin;
    bool myState;
};
\end{cppcode}

\noindent Using the constructor above, we can create a \code{Led} instance and pass it as a copy
(\emph{pass-by-value}) to a function \code{ledPrint}, as shown below. The output is printed using
\code{std::printf()} from \code{<cstdio>}:

\begin{cppcode}
void ledPrint(const Led led) noexcept
{
    const char* state{led.isEnabled() ? "on" : "off"};
    std::printf("The LED is %s!\n", state);
}

int main()
{
    Led led1{8U};
    led1.setEnabled(true);
    ledPrint(led1);
    return 0;
}
\end{cppcode}

\noindent The output appears as shown below:

\begin{console}
The LED is on!
\end{console}

\noindent We can also create a \code{Led} instance through copy initialization, as shown below. Note
that in this case we have not created a \code{Led} instance --- instead, a pin number is passed (by
mistake). This argument is implicitly converted into a \code{Led} instance:

\begin{cppcode}
void ledPrint(const Led led) noexcept
{
    const char* state{led.isEnabled() ? "on" : "off"};
    std::printf("The LED is %s!\n", state);
}

int main()
{
    ledPrint(8U);
    return 0;
}
\end{cppcode}

\noindent The output appears as shown below:

\begin{console}
The LED is off!
\end{console}

\note In this example, \code{led} is passed as a copy (pass-by-value) to clearly demonstrate how an
implicit conversion can occur through the constructor. For larger classes, a const reference is
often used instead to avoid unnecessary copies.

What happens in this case is that:
\begin{itemize}
  \item The compiler implicitly constructs a \code{Led} instance given from an unsigned integer used
    as the pin number.
  \item The compiler expected a \code{Led} instance, but since it received an unsigned integer
    instead, which matches the class constructor, it implicitly invoked this constructor and
    thereby created a \code{Led} instance.
  \item This might be seen as efficient or elegant, but it can also lead to hard-to-detect bugs and
    is therefore not recommended.
  \item It is therefore good practice to ensure that implicit conversion of \code{Led} instances is
    not allowed. This can be achieved by adding the keyword \code{explicit} before the
    constructor, as shown below:
\end{itemize}

\begin{cppcode}
#include <cstdint>

class Led final
{
public:
    explicit Led(const std::uint8_t pin) noexcept
        : myPin{pin}
        , myState{false}
    {}

    void setEnabled(const bool state) noexcept { myState = state; }
    [[nodiscard]] bool isEnabled() const noexcept { return myState; }

private:
    const std::uint8_t myPin;
    bool myState;
};
\end{cppcode}

\noindent After marking the constructor \code{explicit}, the previous code example would have
produced a compilation error.

\subsection{\code{static} members}
\label{c2:sec:static}
So far, every method and member variable we have seen belongs to a specific object: \code{myPin}
and \code{myState} each have a different value for every \code{Gpio} instance, and \code{read()}
operates on one particular instance's state. The keyword \code{static} lets a member instead belong
to the \emph{class itself}, shared by every instance rather than duplicated per object.

In this section we continue using the class \code{Gpio}.

\subsubsection{Static methods}
A static method is called through the class name, using the scope resolution operator \code{::},
rather than through an object:

\begin{cppcode}
Gpio::instanceCount();
\end{cppcode}

\noindent Because a static method has no associated object, it has no \code{this} pointer and
therefore cannot access non-static (instance) members such as \code{myState}. It can only access
other \code{static} members.

We add a static method \code{instanceCount()} that reports how many \code{Gpio} instances currently
exist:

\begin{cppcode}
/**
 * @brief Get the number of currently active GPIO instances.
 *
 * @return Number of active GPIO instances.
 */
[[nodiscard]] static std::uint8_t instanceCount() noexcept { return ourInstanceCount; }
\end{cppcode}

\note A static method can technically also be called through an object, for example
\code{led1.instanceCount()}, but this is discouraged; it misleadingly suggests that the result
depends on \code{led1}, when it does not.

\subsubsection{Static data members}
A static data member is shared by all instances of the class, instead of each object storing its
own copy. We use the prefix \code{our} for private static members, the same way \code{my} is used
for private instance members; both exist to avoid name collisions with methods.

\textbf{\code{static constexpr}} is used for a compile-time constant shared by the class, for
example the maximum number of \code{Gpio} instances supported by the target MCU:

\begin{cppcode}
/** Maximum number of GPIO instances supported by this MCU. */
static constexpr std::uint8_t MaxInstances{40U};
\end{cppcode}

\note Unlike regular \code{camelCase} members and methods, a \code{static constexpr} member is
named starting with a capital letter, for example \code{MaxInstances} rather than
\code{maxInstances}. Some C++ style guides instead use a \code{k} prefix for this
(\code{kMaxInstances}); this book uses a leading capital letter.

Since C++17, a \code{static constexpr} data member is implicitly \code{inline}. This means
\code{MaxInstances} needs no separate definition in a source file. Unlike the plain \code{static}
member shown next, it can be fully declared and initialized directly inside the class.

A \textbf{plain \code{static}} member is used instead when the shared value must change at runtime,
for example a counter tracking how many \code{Gpio} instances currently exist:

\begin{cppcode}
private:
    static std::uint8_t ourInstanceCount;
\end{cppcode}

\noindent Unlike \code{static constexpr}, a plain \code{static} member must be defined exactly once
in a source file, or the linker fails:

\begin{cppcode}
std::uint8_t Gpio::ourInstanceCount{0U};
\end{cppcode}

\noindent We update the constructor and destructor to keep \code{ourInstanceCount} accurate:

\begin{cppcode}
Gpio::Gpio(const std::uint8_t pin, const Direction direction, const bool initialState) noexcept
    : myPin{pin}
    , myDirection{direction}
    , myState{initialState}
{
    ++ourInstanceCount;
}

Gpio::~Gpio() noexcept
{
    --ourInstanceCount;
}
\end{cppcode}

\note This is more than a bookkeeping exercise. Recall from \secref{c2:sec:copymove1} that a
\code{Gpio} object represents a unique physical pin --- \code{ourInstanceCount} combined with
\code{MaxInstances} gives the constructor a way to actually enforce that limit, refusing to create
more \code{Gpio} instances than the MCU has physical pins for.

We can now query the shared instance count without needing any particular \code{Gpio} object:

\begin{cppcode}
int main()
{
    Gpio led{13U, Direction::Output};
    std::printf("Active GPIO instances: %u\n", Gpio::instanceCount());
    return 0;
}
\end{cppcode}

\begin{console}
Active GPIO instances: 1
\end{console}

\subsection{Copy and move operations}
\label{c2:sec:copymove2}
In addition to regular constructors, C++ also provides special member functions for copying and
moving. These are very important in modern C++, and it is useful to become familiar with them when
learning about classes.

The most common ones are:
\begin{itemize}
  \item \textbf{Copy constructor}: creates a new object as a copy of another object.
  \item \textbf{Copy assignment operator}: assigns an already existing object from another object.
  \item \textbf{Move constructor}: creates a new object by moving resources from another object.
  \item \textbf{Move assignment operator}: assigns an already existing object by moving resources
    from another object.
\end{itemize}
In this section we continue using the class \code{Gpio}. This allows the same class to be used
throughout the chapter while demonstrating the syntax for copy and move operations.

The following implementations are mainly shown to illustrate the syntax. For a class such as
\code{Gpio}, it is usually better to delete these operations entirely.

We must keep in mind that the member variables \code{myPin} and \code{myDirection} are declared as
\code{const}. This means that these values can only be set when the object is created and cannot be
changed later. This affects how copy and move operations can be implemented.

\subsubsection{Copy constructor}
The copy constructor is used to create a new object as a copy of another object.

\begin{cppcode}
// Create led1.
Gpio led1{13U, Direction::Output, false};

// Create copy of led1.
Gpio led2{led1};
\end{cppcode}

\noindent An explicit copy constructor can be implemented as follows:

\begin{cppcode}
/**
 * @brief Copy constructor.
 *
 * @param[in] other GPIO instance to copy.
 */
Gpio(const Gpio& other) noexcept
    : myPin{other.myPin}
    , myDirection{other.myDirection}
    , myState{other.myState}
{}
\end{cppcode}

\noindent Here we initialize the member variables of the new instance with the values from the
object \code{other}.

\subsubsection{Move constructor}
The move constructor is used to create a new object by moving resources from another object.

For a simple class such as \code{Gpio}, where we only store small built-in types, there are actually
no heavy resources to move. Despite this, it is still useful to demonstrate the syntax.

\begin{cppcode}
#include <utility>

// Create led1.
Gpio led1{13U, Direction::Output};

// Move led1 into led2.
Gpio led2{std::move(led1)};
\end{cppcode}

\note To invoke the move constructor, the function \code{std::move()} from \code{<utility>} is
used.

An explicit move constructor can be implemented as follows:

\begin{cppcode}
/**
 * @brief Move constructor.
 *
 * @param[in, out] other GPIO instance to move from.
 */
Gpio(Gpio&& other) noexcept
    : myPin{other.myPin}
    , myDirection{other.myDirection}
    , myState{other.myState}
{
    // Clear resources allocated by `other`
    other.myState = false;
}
\end{cppcode}

\note[Notes]
\begin{itemize}
  \item \code{other} is an \textbf{rvalue reference}, written as \code{Gpio&&}. This allows the
    object's resources to be moved instead of copied.
  \item We cannot modify \code{other.myPin} or \code{other.myDirection}, because these member
    variables are declared as \code{const}.
  \item For this class, moving is in practice identical to copying, since no resources are owned
    dynamically.
  \item If \code{other} had contained other resources, such as dynamically allocated memory,
    ownership of these would have been transferred to the new object, and \code{other} would have
    been left without them (for example with its pointer set to \code{nullptr}).
\end{itemize}

\subsubsection{Copy assignment operator}
A copy assignment operator is used when an already existing object should be assigned from another
object.

\begin{cppcode}
// Create led1 and led2.
Gpio led1{13U, Direction::Output};
Gpio led2{12U, Direction::Output};

// Make led2 a copy of led1.
led2 = led1;
\end{cppcode}

\noindent For the class \code{Gpio}, however, this is problematic. During assignment, an already
existing object must be overwritten, but \code{myPin} and \code{myDirection} are constant and may
not be changed after initialization.

An otherwise simplified operator could have looked like this:

\begin{cppcode}
/**
 * @brief Copy assignment operator.
 *
 * @param[in] other GPIO instance to copy from.
 *
 * @return Reference to this instance.
 */
Gpio& operator=(const Gpio& other) noexcept
{
    if (this != &other)
    {
        myState = other.myState;
    }
    return *this;
}
\end{cppcode}

\noindent Note that we \textbf{cannot} write:

\begin{cppcode}
myPin       = other.myPin;
myDirection = other.myDirection;
\end{cppcode}

\noindent since these are constant.

\subsubsection{Move assignment operator}
Similarly, there is a move assignment operator, which is used when an object is assigned from a
temporary object or via \code{std::move} from \code{<utility>}:

\begin{cppcode}
#include <utility>

Gpio led1{13U, Direction::Output};
Gpio led2{12U, Direction::Output};

led2 = std::move(led1);
\end{cppcode}

\noindent A simplified implementation could look like this:

\begin{cppcode}
/**
 * @brief Move assignment operator.
 *
 * @param[in, out] other GPIO instance to move from.
 *
 * @return Reference to this instance.
 */
Gpio& operator=(Gpio&& other) noexcept
{
    if (this != &other)
    {
        myState       = other.myState;
        other.myState = false;
    }
    return *this;
}
\end{cppcode}

\noindent Here as well, however, we cannot move \code{myPin} or \code{myDirection}, since these are
constant.

\subsection{Recommendation}
\label{c2:sec:recommendation}
In many real embedded systems, copying and moving are often completely disallowed, since an object
frequently represents a unique hardware resource.

The deleted copy and move operations are normally placed at the bottom of the public section:

\begin{cppcode}
Gpio(const Gpio&)            = delete; // No copy constructor.
Gpio(Gpio&&)                 = delete; // No move constructor.
Gpio& operator=(const Gpio&) = delete; // No copy assignment.
Gpio& operator=(Gpio&&)      = delete; // No move assignment.
\end{cppcode}

\noindent In this way, it is ensured that each \code{Gpio} object represents exactly one physical
GPIO pin and cannot be duplicated by mistake.
